\documentclass{sprawozdanie-agh}

\usepackage[utf8]{inputenc}
\usepackage{matlab-prettifier}
\usepackage{siunitx}
\usepackage{booktabs,caption,ragged2e}

\makeatletter

% ========= Początek Dokumentu ========= 

\begin{document}
% ========= Strona tytułowa ========= 
\przedmiot{Teoria Sterowania II: Ćwiczenia Laboratoryjne}
\tytul{Wytyczne dla Ćwiczenia 3}
\podtytul{Metoda Lapunowa}
\kierunek{Automatyke i~Robotyka}
\autor{Dariusz Cieślar, Maciej Różewicz}
\data{Kraków, 3 marca 2026}

\stronatytulowa{}

% ========= Treść Właściwa =========
\section{Cel ćwiczenia}
Celem ćwiczenia jest zapoznanie się z~metodami badania stabilności układów nieliniowych.
W czasie zajęć przedstawione i~przećwiczone zostaną dwie metody badania stabilności:
\begin{itemize}
    \item pośrednia metoda Lapunowa - prosta metoda, dająca zadowalające wyniki w~wielu praktycznych przypadkach, jednak mająca pewne ograniczenia,
    \item bezpośrednia metoda Lapunowa - bardzo uniwersalna metoda, pozwalająca zbadać stabilność szerokiej gamy układów (zarówno stacjonarnych, jak i~niestacjonarnych).
\end{itemize}
Ponadto po sprawdzeniu stabilności punktów równowagi układu badany i~estymowany będzie 
obszar przyciągania asymptotycznego za pomocą twierdzenia La Salle'a.
Dodatkowo sprawdzana będzie stabilność nie tylko punktów równowagi, ale i~cykli granicznych.

%%%%%%%%%%%%%%%%%%%%%%%%%%%%%%%%%%%%%%%%%%%%%%%%%%%%%%%%%%%% 
\section{Przebieg ćwiczenia}
    %=======================================================
    \subsection{Pośrednia metoda Lapunowa}
    Dla każdego z~podanych systemów opisanych równaniami (\ref{eq:posrednia_1}), (\ref{eq:posrednia_2}), (\ref{eq:posrednia_3}),  (\ref{eq:posrednia_4}),  (\ref{eq:posrednia_5}):
    \begin{itemize}
        \item zamodelować układ równań w~Simulinku (lub użyć skryptu do jego rozwiązania),
        \item wyznaczyć portrety fazowe,
        \item znaleźć punkty równowagi,
        \item sprawdzić, czy można zastosować pośrednią metodę Lapunowa do badania stabilności punktów równowagi,
        \item zbadać stabilność wszystkich punktów równowagi,
        \item zbadać doświadczalnie obszar atrakcji asymptotycznie stabilnych punktów równowagi - przedstawić go na tle portretu fazowego.
    \end{itemize}
        %---------------------------------------------------
        \subsubsection{System 1}
        Rozważamy układ mechaniczny masa-sprężyna ze sprężyną o~nieliniowej charakterystyce.
        Układ ten opisany jest równaniem (\ref{eq:posrednia_1}).
        \begin{equation}\label{eq:posrednia_1}
            \ddot{x} + b\dot{x} + cx + dx^{3} = 0, \, b,c > 0,\, |c|<|d|.
        \end{equation}
        Gdzie:
        \begin{itemize}
            \item $x$ - odchylenie od punktu równowagi,
            \item $b$ - współczynnik tarcia,
            \item $c$ i~$d$ - parametry sprężyny.
        \end{itemize}
        Proszę rozważyć przypadek dla $d>0$ i~$d<0$.

        %---------------------------------------------------
        \subsubsection{System 2}
        Rozważamy wahadło z~tłumieniem, opisywane równaniem:
        \begin{equation}\label{eq:posrednia_2}
            \ddot{x} + \frac{c}{lm}\dot{x} + \frac{g}{l}\sin{x} = 0
        \end{equation}
        gdzie:
        \begin{itemize}
            \item $x$ - kąt wychylenia wahadła z~punktu równowagi trwałej,
            \item $g$ - przyspieszenie ziemskie,
            \item $m$ - masa wahadła,
            \item $l$ - długość wahadła.
        \end{itemize}

        %---------------------------------------------------
        \subsubsection{System 3}
        Rozważamy układ równań Van der Pola:
        \begin{equation}
            \label{eq:posrednia_3}
            \left\{
            \begin{array}{l}
                 \dot{x}_{1} = x_{2} - x_{1}^{3} - ax_{1}\\
                 \dot{x}_{2} = -x_{1}
            \end{array}
            \right\}.
        \end{equation}
        Proszę zbadać zachowanie układu dla różnych wartości $a$.

        %---------------------------------------------------
        \subsubsection{System 4}
        Rozważamy układ równań:
        \begin{equation}
            \label{eq:posrednia_4}
            \left\{
            \begin{array}{l}
                 \dot{x}_{1} = -x_{1} + \sin{x_{2}}\\
                 \dot{x}_{2} = -x_{2}
            \end{array}
            \right\}.
        \end{equation}

        %---------------------------------------------------
        \subsubsection{System 5}
        Rozważamy układ równań:
        \begin{equation}
            \label{eq:posrednia_5}
            \left\{
            \begin{array}{l}
                 \dot{x}_{1} = 2x_{1} - x_{2} - x_{1}^{3}\\
                 \dot{x}_{2} = 4x_{1} - 2x_{2}
            \end{array}
            \right\}.
        \end{equation}

    %=======================================================
    \subsection{Bezpośrednia metoda Lapunowa i~twierdzenie La Salle'a}
    Dla poniższych układów (\ref{eq:zadania_przyklad_1}), (\ref{eq:zadania_przyklad_2}), (\ref{eq:zadania_przyklad_3}), (\ref{eq:zadania_przyklad_4}):
    \begin{itemize}
        \item zamodelować układ równań w~Simulinku (lub użyć skryptu do jego rozwiązania),
        \item wyznaczyć portrety fazowe,
        \item znaleźć punkty równowagi,
        \item zbadać stabilność punktów równowagi, stosując podane funkcje Lapunowa,
        \item dla każdej z~podanych funkcji Lapunowa znaleźć analitycznie estymatę obszaru atrakcji,
        \item symulacyjnie określić obszar atrakcji,
        \item porównać obszary odnalezione analitycznie i~symulacyjnie - przedstawić obydwa na tle portretu fazowego,
        \item porównać estymaty obszaru atrakcji uzyskane z~użyciem różnych funkcji Lapunowa - przedstawić obydwa na jednym wykresie na tle portretu fazowego.
    \end{itemize}
        %---------------------------------------------------
        \subsubsection{System 1}
        Rozważamy układ:
        \begin{equation}
            \label{eq:zadania_przyklad_1}
            \left\{
            \begin{array}{l}
                 \dot{x}_{1} = -x_{1} + 2x_{1}^{2}x_{2}\\
                 \dot{x}_{2} = -x_{2}
            \end{array}
            \right\}.
        \end{equation}
        Dla tego układu rozpatrujemy dwóch kandydatów na funkcję Lapunowa:
        \begin{itemize}[(a)]
            \item $V_{1}^{1}(x) = \frac{1}{2}x_{1}^{2} + x_{2}^{2}$,
            \item $V_{1}^{2}(x) = \frac{x_{1}^{2}}{1 - x_{1}x_{2}} + x_{2}^{2}$.
        \end{itemize}

        %---------------------------------------------------
        \subsubsection{System 2}
        Rozważamy układ:
        \begin{equation}
            \label{eq:zadania_przyklad_2}
            \left\{
            \begin{array}{l}
                 \dot{x}_{1} = x_{2} - x_{1} + x_{1}^{3}\\
                 \dot{x}_{2} = -x_{1}
            \end{array}
            \right\}.
        \end{equation}
        Dla tego układu rozpatrujemy dwóch kandydatów na funkcję Lapunowa:
        \begin{itemize}[(a)]
            \item $V_{2}^{1}(x) = \frac{1}{2} \left( x_{1}^{2} + x_{2}^{2} \right)$,
            \item $V_{2}^{2}(x) = \frac{1}{2} \left( 2x_{1}^{2} + 2x_{1}x_{2} + 3x_{2}^{2} \right)$.
        \end{itemize}

        %---------------------------------------------------
        \subsubsection{System 3}
        Rozważamy układ (zauważmy, że nie rozważamy $x\in \mathbb{R}^{2}$, ale $x_{1} > -1$):
        \begin{equation}
            \label{eq:zadania_przyklad_3}
            \left\{
            \begin{array}{l}
                 \dot{x}_{1} = x_{2}\\
                 \dot{x}_{2} = -\frac{x_{2}}{1 + x_{1}} - \frac{x_{1}}{1 + x_{1}}
            \end{array}
            \right\}.
        \end{equation}
        Dla tego układu rozpatrujemy dwóch kandydatów na funkcję Lapunowa:
        \begin{itemize}[(a)]
            \item $V_{3}^{1}(x) = \frac{1}{2}x_{2}^{2} + \int_{0}^{x_{1}}{\frac{z}{1+z}dz}$,
            \item $V_{3}^{2}(x) = x_{1}^{2} + x_{1}x_{2} + x_{2}^{2}$.
        \end{itemize}

        \subsubsection{System 4}
        Rozważamy układ (szukamy cyklu granicznego):
        \begin{equation}
            \label{eq:zadania_przyklad_4}
            \left\{
            \begin{array}{l}
                 \dot{x}_{1} = x_{2} - x_{1}^7 \left( x_{1}^{4} + 2x_{2}^{2} - 10 \right)\\
                 \dot{x}_{2} = -x_{1}^{3} - 3x_{2}^{5} \left( x_{1}^{4} + 2x_{2}^{2} - 10 \right)
            \end{array}
            \right\}.
        \end{equation}
        Dla tego układu rozpatrujemy następującego kandydata na funkcję Lapunowa:
        \begin{itemize}[(a)]
            \item $V_{4}^{1}(x) = \left( x_{1}^{4} + 2x_{2}^{2} - 10 \right)^{2}$.
        \end{itemize}



\end{document}
